%% kapa: exam=1 round=31 year=2020 period=1 form=A subject=civil_law subject_ko=민법 seq=21 number=21
%% source: Q-Net 감정평가사 기출문제 (cst00309) artlSeq=5207219 file=extracted/1008/5207219/A형/A형/2020년도 제31회 감정평가사 1차 시험_1교시_A형.pdf
%% conversion: kordoc 4.24.51 -> tools/gen1.py (auto); review=proofread-sonnet; flags=none
\begin{question}{21}
물권에 관한 설명으로 옳지 않은 것은? (다툼이 있으면 판례에 따름)
\begin{choices}
  \item 특별한 사정이 없으면, 물건의 일부는 물권의 객체가 될 수 없다.
  \item 권원 없이 타인의 토지에 심은 수목은 독립한 물권의 객체가 될 수 없다.
  \item 종류, 장소 또는 수량지정 등의 방법으로 특정할 수 있으면 수량이 변동하는 동산의 집합도 하나의 물권의 객체가 될 수 있다.
  \item 소유권을 비롯한 물권은 소멸시효의 적용을 받지 않는다.
  \item 소유권을 상실한 전(前)소유자는 물권적 청구권을 행사할 수 없다.
\end{choices}
\end{question}
\answer{4}
